\ExampleChapter{Abstract}{Abstract}
Multivariate flight records provide a useful setting for illustrating time-series
modeling and anomaly analysis. This document presents a reproducible example
based entirely on synthetic signals. It is intended to demonstrate the
GZUThesis workflow rather than to introduce a new detection algorithm or
report results from real aircraft.

Six channels are generated from periodic components, cross-channel terms, and
random noise. Drift, vibration, and bias anomalies are injected into held-out
synthetic flights. Complete flights are assigned to training, calibration, or
testing; no sliding window crosses a flight boundary. A ridge predictor is
fitted on standardized training observations. Its one-step residuals are
aggregated across channels, optionally smoothed with an exponentially weighted
moving average, and compared with a threshold estimated from normal calibration
scores. Persistence and moving-average predictors provide transparent baselines.

For the bundled fixed-seed experiment, ridge prediction with smoothing obtains
pointwise precision, recall, and F1 scores of \DemoPrecision\%, \DemoRecall\%,
and \DemoFone\%, respectively. These figures characterize only the generated
example and do not establish performance on real UAV failures. Results are
reported without point adjustment. The discussion separates prediction,
post-processing, calibration, and evaluation so that changes in one component
are not mistaken for improvements in another. The accompanying project includes
source chapters, figures, tables, a shared bibliography, and executable scripts,
allowing both the numerical outputs and the document to be rebuilt.

\par\vspace{1em}
\noindent\textbf{Keywords:} unmanned aerial vehicle; time series; anomaly detection;
prediction residual; reproducible typesetting
